% sec-10-relevant-information.tex - Relevant Information (final stage)
\section{Relevant Information}

This section provides the additional relevant information that the U.S. Food and
Drug Administration may request under 21 CFR \S 312.23(a)(11), together with the
method-level grounding and verification disclosures that apply because this
Investigational New Drug application (IND v1.0) was assembled by an artificial
intelligence build operating over the repository at version v4.3.0. The
disclosures are placed here, at the close of the package, so that a reviewer can
read the entire submission first, on its clinical, statistical, and regulatory
merits, and then evaluate exactly how the document was constructed: how each
figure and number was grounded to a named source file, how each figure was cited
where its content is discussed, how the build is monitored in real time on a
public version-control system, how the repository and file naming mirror
document structure, and where residual limitations of method lie. 

Nothing
in this section substitutes for the clinical, statistical, and regulatory
substance of the preceding sections; rather, it states the controls that give a
reviewer confidence that the substance is internally consistent across the
package, traceable to named source files, and reproducible from those sources by a
third party. The verification posture described here is aligned with the
verification-before-generation discipline that the program has advanced for
physical artificial intelligence in oncology trials, the same discipline codified
in the H.R. 9510 framing on which the device-directed component of the trial rests
\cite{congress9510, clintrustpai}, and it is consistent with the
phased-development expectations that govern an initial Phase 1 submission
\cite{phase1guidance}. The disclosures should be read as a transparency artifact
offered in the spirit of 21 CFR \S 312.23(a)(11), not as a request for any special
regulatory accommodation: every generated document, this IND included, remains
subject to the same human medical, statistical, and regulatory review as a
conventionally authored submission.

The structure of the section is as follows. We first state the method-level
grounding and verification discipline and present, in summary and then in detail,
the five name-matching and grounding verifications that the build applies to every
artifact it emits (Figure 21 and Table 11.1). We then describe the finer-grained
figure-grounding discipline that keeps each rendered figure provably identical to
its Mermaid source (Figure 22 and Table 11.2). We next document the real-time
auto-commit and auto-pull-request monitorability of the build, the repository,
directory, and file name matching that makes the package self-describing (Table
11.3), and the principal limitations of the method together with the mitigation
applied to each (Table 11.4). A short worked example ties the abstractions to a
single concrete trial number (Table 11.5), and a closing compilation statement
records how a reviewer can reproduce the compiled PDF directly from the repository.

\subsection{Method-Level Grounding and Verification}

The IND was produced by a single-prompt, multi-file build that proceeds through
four stages, mermaid to draft to full to final, with the present final stage
expanding each full-stage section toward the package length target and applying the
senior-author formatting pass that a PDF compiler is required to verify. To keep
that build verifiable rather than merely fluent, the method applies five
name-matching and grounding verifications to every artifact it emits. The five
verifications are summarized in Figure 21 and detailed, with a concrete IND example
for each, in Table 11.1. They are not abstract quality slogans; each one is a check
that a reviewer can independently reproduce by inspecting the public repository
tree, the rendered document, and the commit and pull-request history. The five
together yield a grounded, monitorable IND in which every figure names its source,
every figure is cited where its section discusses it, every file is visible on
GitHub as it is written, repository and directory names mirror the document
structure, and section file names map one to one onto regulatory sections
they contain.

This discipline matters for a regulatory submission in a way that distinguishes it
from generic large-language-model authoring. A submission to FDA is read for
internal consistency: a dose level stated in the General Investigational Plan must
match the dose level in the protocol synopsis, in the dose-escalation figure, and
in the safety narrative; a force cap stated in the device description must match
the cap in the safety-reporting triggers. Ungrounded generation is fluent but does
not guarantee that the same number appears identically wherever it is reused, and a
single drifted number in a Phase 1 oncology submission can change the reviewer's
understanding of the risk. The grounding discipline forecloses that drift by
binding every quantitative claim a figure carries to a named source file from which
the surrounding prose draws the same value, so that the figure and the text cannot
diverge \cite{oncotrialllm, pdacdigtwinp}.

\needspace{8\baselineskip}
\begin{mermaidfig}
\node[mmdec] (k1) at (0,-0.15) {1 Grounding};
\node[mmdec] (k2) at (0,-2.0) {2 Figures match context};
\node[mmdec] (k3) at (0,-4.0) {3 GitHub real-time};
\node[mmdec] (k4) at (0,-6.0) {4 Repo / dir names match};
\node[mmdec] (k5) at (0,-8.0) {5 File names match};
\node[mmin] (e1) at (5.0,-0.15) {22 grayscale Mermaid figures, each sourced};
\node[mmin] (e2) at (5.0,-2.0) {each figure cited where its \S section discusses it};
\node[mmin] (e3) at (5.0,-4.0) {one commit per file, auto-pushed to one PR};
\node[mmin] (e4) at (5.0,-6.0) {trial-ind/ matches the repository tree};
\node[mmin] (e5) at (5.0,-8.0) {sec-06-cmc.tex matches the CMC section};
\node[mmgoal] (g) at (10.4,-4.0) {Grounded,\\monitorable IND};
\draw[mmedge] (k1) -- (e1);
\draw[mmedge] (k2) -- (e2);
\draw[mmedge] (k3) -- (e3);
\draw[mmedge] (k4) -- (e4);
\draw[mmedge] (k5) -- (e5);
\draw[mmedge] (e1) to[out=0,in=150,looseness=0.6] (g.west);
\draw[mmedge] (e2) to[out=0,in=165,looseness=0.6] (g.west);
\draw[mmedge] (e3) -- (g.west);
\draw[mmedge] (e4) to[out=0,in=195,looseness=0.6] (g.west);
\draw[mmedge] (e5) to[out=0,in=210,looseness=0.6] (g.west);
\end{mermaidfig}
\vspace*{-0.6cm}
\figcaption{Figure 21. Grounding (every one of the 22 grayscale figures names its
source files); figures match context (each is cited where its section discusses
it); GitHub real-time (one commit per file, auto-pushed to one continuously
updated pull request); repository and directory names match (the trial-ind tree
mirrors the document structure); and file names match (sec-06-cmc.tex maps to the
Chemistry, Manufacturing and Control section).}

\subsubsection{The Five Verifications in Detail}

\textit{Verification 1, grounding.} Every quantitative claim that a figure carries
is traceable to a named source file rather than invented at generation time. The
22 grayscale figures of this IND each declare, in their Mermaid source under
\texttt{trial-ind/mermaid/}, the source files from which their nodes, edge labels,
and numbers were drawn, and those numbers are the same document counts, dose
levels, force caps, sampling rates, and review clocks reused in the surrounding
prose. For example, the dose levels of 160 mg, 220 mg, and 300 mg once daily, the
28 day dose-limiting-toxicity window, and the maximum enrollment of 18 treated
subjects appear identically in the dose-escalation figure, in the General
Investigational Plan, and in the protocol from which they originate. The same
binding holds for the device parameters that recur across the safety sections:
the eight robotic arms, the 56 total degrees of freedom (8 arms times 7 each), the
640 sensor channels (80 per arm), the 100 kHz force sampling and 10 kHz sampling
for other channels, the per-arm force cap of 3 N and the cumulative cap of 18 N,
the 10 kHz safety heartbeat with its 100 microsecond watchdog, the cross-arm
emergency stop at 3 ms to 0.0 N, and the hardware emergency stop at 500 ms under
21 CFR \S 312.404 all appear with identical values wherever they are reused. This
binds the figures to the trial facts and prevents a figure from drifting away from
the text it illustrates \cite{oncotrialllm, pdacdigtwinp}.

\textit{Verification 2, figures match context.} Each figure is cited in the
sentence or paragraph that discusses it, in the section where that discussion
belongs, so that no figure floats unreferenced and no claim in the prose lacks its
illustrating figure. The daraxonrasib mechanism-of-action figure is cited in the
introductory material that names the drug and its pharmacological class, the 3+3
dose-escalation automaton is cited in the General Investigational Plan, the
chemistry information architecture is cited in the Chemistry, Manufacturing and
Control section, the perioperative pause-and-restart advisory is cited where the
restart-day decision is described, and the two figures of the present section are
cited here, as Figure 21 and Figure 22. A reviewer can confirm this by following
any figure number from its caption back to the citing sentence, and the converse:
every figure number that appears in prose resolves to one figure.

\textit{Verification 3, GitHub real time.} Each distinguishable file is committed
and pushed at the moment it is written, and those commits accumulate on a single
continuously updated pull request, so that the IND is observable on GitHub as it is
constructed rather than only after completion. This gives the sponsor-investigator,
and in principle an auditor, a real-time view of provenance: who or what wrote each
file, in what order, and against which commit. The auto-commit and
auto-pull-request behavior is the documentary analogue of the audit-trail and
monitorability expectations that govern the device-directed component of the trial,
under which the device preserves a window of records from 24 hours before to 72
hours after each physical artificial intelligence adverse event under 21 CFR part
11 \cite{paiautospons, paisitedocpk}.

\textit{Verification 4, repository and directory names match.} The directory tree
named in the prose is the directory tree on disk. References to
\texttt{trial-ind/mermaid/}, \texttt{trial-ind/full-ind/sections/},
\texttt{trial-ind/final-ind/sections/}, and \texttt{trial-ind/inputs/} denote the
actual directories of the repository at version v4.3.0, so a reader can navigate
from any path mentioned in the document straight to the corresponding location in
the public tree without translation. The four stage directories, \texttt{mermaid/},
\texttt{draft-ind/}, \texttt{full-ind/}, and \texttt{final-ind/}, mirror the four
build stages in the same left-to-right order in which the prose describes them.

\textit{Verification 5, file names match.} The section file names map one to one
onto the regulatory sections they contain. sec-06-cmc.tex is the
Chemistry, Manufacturing and Control section, sec-07-pharmacology-toxicology.tex
is the Pharmacology and Toxicology section, sec-10-relevant-information.tex
is the present section, and fig-09-five-grounding-verifications.md is the
Mermaid source for Figure 21. Naming convention makes package
self-describing and removes ambiguity for which file holds which content, so a
reviewer can confirm that no section is missing and no file is misnamed simply by
listing directory.

\needspace{10\baselineskip}
\noindent\begin{tabularx}{\textwidth}{L{0.5cm} L{4.2cm} Y}
\toprule
\textbf{\#} & \textbf{Verification check} & \textbf{Concrete IND example} \\
\midrule
1 & Grounding via Mermaid figures & The 22 grayscale figures in \texttt{trial-ind/mermaid/} each name their source files and carry the dose levels (160, 220, 300 mg QD), the 28 day DLT window, the force caps (3 N per arm, 18 N cumulative), and the 30 day FDA clock reused verbatim in the prose \\
2 & Figures match IND context & Each figure is cited in the sentence that discusses it, in the section where that discussion belongs (mechanism in the Introduction, 3+3 automaton in the Plan, CMC architecture in the CMC section, these two figures here) \\
3 & AI files viewable on GitHub in real time & One commit per file is pushed at generation, accumulating on a single continuously updated pull request, so each artifact appears on the branch as it is made \\
4 & Repository and directory names match & The prose refers to \texttt{trial-ind/mermaid} and \texttt{full-ind/sections}, which are the directories on disk in repository v4.3.0 \\
5 & File names match repository file names & A cited \texttt{sec-06-cmc.tex} is the CMC section; \texttt{fig-09-five-grounding-verifications.md} is the Mermaid source for Figure 21 \\
\bottomrule
\end{tabularx}
\vspace*{-0.35cm}
\figcaption{Table 11.1. Five name-matching and grounding verifications, each
with an example drawn from this IND. Checks are independently
reproducible from public repository tree, rendered document, and commit
and pull-request.}

\subsubsection{Why Grounding Is the Load-Bearing Control}

Of the five verifications, grounding is the one on which the regulatory value of
the package most depends, and it is worth stating precisely what it does and does
not guarantee. Grounding guarantees that a number rendered in a figure or asserted
in prose is the same number that appears in the named source file from which it was
drawn; it does not, by itself, guarantee that the source number is clinically
correct. Clinical correctness is established by the human medical, statistical, and
regulatory review that every generated document must pass before use, and by the
underlying protocol, investigator brochure, and chemistry record that the build
ingests as curated inputs. What grounding contributes is the elimination of a
specific and common failure mode of language-model authoring, namely the silent
substitution of a plausible but unsourced value. Because the dose levels, the DLT
window, the enrollment cap, the device force caps, and the safety-reporting clocks
are each carried from a single source through every figure and every paragraph that
reuses them, a reviewer who verifies one occurrence has effectively verified all of
them, and any divergence would be visible as a mismatch against the named source.
This is the property that lets a reviewer audit the package efficiently rather than
re-checking every restatement of every number independently
\cite{oncotrialllm, llmcomp16fda}.

The grounding discipline is reinforced by the fixed grayscale encoding of node
roles in the figures. Because the eight-tone ramp is defined once in
\texttt{indstyle.sty} and every figure draws its node fills from that ramp, the
tonal meaning of a node, goal, process, accent, step, input, context, decision, or
dark, is identical across all 22 figures and across all four build stages. A
decision diamond in the 3+3 automaton carries the same tone as a decision diamond
in the perioperative advisory or in the figure-grounding diagram of the present
section, so a reviewer learns the visual vocabulary once and reads every figure
with it. This consistency is itself a grounding property: it ties the visual
encoding to a single source of truth in the style file rather than to per-figure
choices that could drift \cite{fdadigtwinpc}.

\subsection{Figure Grounding: Mermaid to TikZ}

A second, finer-grained discipline governs the figures themselves: every figure in
this IND is authored once as a grayscale Mermaid source in
\texttt{trial-ind/mermaid/} and then reproduced as a TikZ figure of identical
complexity in the LaTeX document, with no simplification of nodes, edges, edge
labels, or quantitative content. The translation pathway and its verify-twice
checkpoint are shown in Figure 22. The Mermaid source defines the nodes, edges,
grayscale palette, and quantitative data; the conversion preserves all of these;
the TikZ \texttt{mermaidfig} environment renders them using the \texttt{mm*} node
styles defined in \texttt{indstyle.sty} (the same eight-tone ramp that maps one to
one from the Mermaid classDef classes); and the result is verified twice, for
text-box and arrow overlaps, for the specified arrow looseness, and for proper box
spacing, before it appears identically in the draft, full, and final stages. A
figure that fails the verify-twice checkpoint is returned to the TikZ step and
corrected; only a figure that passes is allowed to render in the document.

\needspace{8\baselineskip}
\begin{mermaidfig}
\node[mmin] (src) at (0,0) {mermaid/fig-NN.md\\(GitHub Mermaid)};
\node[mmproc] (con) at (4.4,0) {nodes, edges, grayscale palette, quantitative data};
\node[mmproc] (tikz) at (8.8,0) {TikZ mermaidfig\\(mm* node styles)};
\node[mmgoal] (out) at (13.2,0) {Identical figure in draft / full / final};
\node[mmdec] (v) at (13.2,-3.2) {Verify twice: no overlaps, arrow looseness, box spacing};
\draw[mmedge] (src) -- (con);
\draw[mmedge] (con) -- (tikz);
\draw[mmedge] (tikz) -- (out);
\draw[mmedge] (out) -- (v);
\draw[mmedge] (v.east) to[out=90,in=-45,looseness=0.6] node[mmlabel,pos=0.55] {pass} (out.east);
\draw[mmedged] (v.north west) to[out=150,in=-60,looseness=0.6] node[mmlabel,pos=0.5] {fail} (tikz.south);
\end{mermaidfig}
\vspace*{-0.65cm}
\figcaption{Figure 22. Figure grounding. Each Mermaid source in \texttt{mermaid/} is
reproduced as a TikZ \texttt{mermaidfig} of the same complexity (the same nodes,
edges, grayscale palette, and quantitative data), then verified twice for text-box
and arrow overlaps, for the specified arrow looseness, and for proper box spacing,
before it appears identically in the draft, full, and final stages. A figure that
fails is returned to the TikZ step; only a figure that passes renders.}

The practical effect of this discipline is that the figure a reviewer sees in the
compiled PDF is provably the same diagram, with the same numbers, as the Mermaid
source committed to the repository. There is no separate figure-rendering pipeline,
no raster export, and no opportunity for a number to be transcribed differently
between the source and the rendered figure. All figures are vector TikZ; the IND
contains no PNG or JPG images and no \texttt{\textbackslash includegraphics} calls.
The grayscale ramp is fixed in the style file, so the tonal encoding of node roles
(goal, process, accent, step, input, context, decision, dark) is identical across
all 22 figures and across the four build stages \cite{fdadigtwinpc}.

\subsubsection{The Verify-Twice Checkpoint}

The verify-twice checkpoint is named for the fact that each figure is inspected on
two independent criteria before it is admitted to the document, and that the
inspection is repeated after any correction so that a fix to one criterion cannot
silently break the other. The first criterion is geometric: no two text boxes
overlap and no arrow crosses through a box. The build enforces explicit spacing
rules that a reviewer can verify against the rendered coordinates, namely that
side-by-side box centers are at least 4.0 cm apart in the horizontal direction,
that stacked box centers are at least 1.7 cm apart in the vertical direction, that
the node text width is 30 mm, and that decision diamonds, which are wider than
rectangles for the same text, are given extra room. The second criterion is the
routing of curved arrows: every curved edge uses a \texttt{to[out,in,looseness]}
specification with a looseness held between 0.4 and 0.7, which keeps the curve
gentle enough to avoid sweeping across a neighboring box while still separating
parallel edges that converge on the same target. Table 11.2 records these geometric
checks together with the threshold applied to each, so that a reviewer can confirm
the layout discipline against the rendered figure rather than taking it on
assertion \cite{fdadigtwinpc, natlpaiplatf}.

\needspace{10\baselineskip}
\noindent\begin{tabularx}{\textwidth}{L{4.6cm} L{3.2cm} Y}
\toprule
\textbf{Geometric check} & \textbf{Threshold} & \textbf{Purpose} \\
\midrule
Side-by-side box centers & $\geq$ 4.0 cm in x & Prevents horizontal text-box overlap at the fixed 30 mm node text width \\
Stacked box centers & $\geq$ 1.7 cm in y & Prevents vertical overlap of stacked boxes and their captions \\
Node text width & 30 mm fixed & Uniform box sizing so spacing thresholds apply identically to every figure \\
Decision diamonds (mmdec) & Extra room reserved & Diamonds are wider than rectangles for the same text and need clearance \\
Curved-arrow looseness & 0.4 to 0.7 & Keeps curves gentle so arrows do not sweep across neighboring boxes \\
Arrow-box intersection & None permitted & No edge may pass through a node it does not connect \\
\bottomrule
\end{tabularx}
\vspace*{-0.35cm}
\figcaption{Table 11.2. The verify-twice geometric checks applied to every figure
before it is admitted to the document, with the threshold and purpose of each. The
checks are repeated after any correction so that a fix to one cannot silently break
another; only a figure that passes both inspections renders in the draft, full, and
final stages.}

Because the checkpoint is geometric and the thresholds are explicit, it is
deterministic in a way that a subjective "does this look right" review is not. Two
reviewers applying the 4.0 cm and 1.7 cm thresholds to the same rendered figure
reach the same verdict, and a figure that passes in the draft stage continues to
pass in the full and final stages because the node coordinates and the style file
are unchanged across stages. This is what allows the program to assert that the
figure is identical from Mermaid source to final PDF: not merely that it looks
similar, but that it satisfies the same enumerated geometric and quantitative
criteria at every stage \cite{fdadigtwinpc}.

\subsection{Real-Time Auto-Commit and Auto-Pull-Request Monitorability}

The build is monitorable in real time. As each section file, style file, figure
source, and supporting file is written, it is committed individually and pushed to
the remote, and the accumulating commits are surfaced on a single open pull request
that updates continuously over the life of the build. This means that the
provenance of the IND is not reconstructed after the fact from a final snapshot;
it is visible while the document is being produced. A reader with access to the
public repository can watch the \texttt{trial-ind/final-ind/} tree fill in section
by section, can inspect the diff for any individual file, and can read the commit
history as an ordered, timestamped record of how the package came together. This
real-time monitorability is the documentary counterpart of the trial's own
audit-trail commitments, under which the device preserves a window of records from
24 hours before to 72 hours after each physical artificial intelligence adverse
event, and it is consistent with the broader transparency posture that the program
has argued is necessary before higher-autonomy systems are used in oncology trials
\cite{congress9510, clintrustpai, natlpaiplatf}. The same single-prompt,
auto-commit, auto-pull-request method has been demonstrated across prior multi-file
builds in the program, including the companion Phase 2 protocol and the
legislative-text repositories, which establishes that the approach is reproducible
rather than specific to this submission \cite{oncotrialllm, paipredict4x}.

It is important to be precise about what this monitorability does and does not
provide in a regulatory sense. It provides a tamper-evident, timestamped,
ordered record of document construction that an auditor can replay commit by
commit; it does not provide, and is not offered as, the formal part 11 audit trail
that governs the trial's clinical and device data, which is a separate system with
its own controls. The two are analogues, not the same system: the build's
version-control history documents how the application was written, while the
trial's part 11 audit trail documents how clinical and device events are recorded
during conduct of the study. Stating this distinction plainly avoids any
implication that authoring a submission with a monitorable build relieves the
sponsor-investigator of the independent obligation to maintain part 11 audit trails
during the trial itself \cite{paiautospons, paisitedocpk}.

\subsection{Repository, Directory, and File Name Matching}

The repository, directory, and file naming is itself a control. The IND lives under
\texttt{trial-ind/} in the public repository, and within it the four build stages
occupy \texttt{mermaid/}, \texttt{draft-ind/}, \texttt{full-ind/}, and
\texttt{final-ind/}, mirroring the mermaid to draft to full to final progression.
Within each LaTeX stage the section files follow a fixed
\texttt{sec-NN-<name>.tex} convention whose number and name match the ReGARDD table
of contents, so that the file that holds the Chemistry, Manufacturing and Control
section is named \texttt{sec-06-cmc.tex}, the file that holds the Pharmacology and
Toxicology section is named \texttt{sec-07-pharmacology-toxicology.tex}, and the
figure that this section reproduces first is sourced from
\texttt{fig-09-five-grounding-verifications.md}. Table 11.3 records the directory
and file-naming correspondences that make the package self-describing, so that a
reviewer can move from any path named in the prose directly to the corresponding
location on disk without translation, and can confirm that no section is missing
and no file is misnamed.

\needspace{10\baselineskip}
\noindent\begin{tabularx}{\textwidth}{L{7.5cm} L{3.5cm} Y}
\toprule
\textbf{Path in repository v4.3.0} & \textbf{Holds} & \textbf{Name-match guarantee} \\
\midrule
\texttt{trial-ind/} & The full IND across four stages & Top-level directory named for the IND it contains \\
\texttt{trial-ind/mermaid/} & The 22 grayscale figure sources & Stage 1 directory; \texttt{fig-NN-<name>.md} files \\
\texttt{trial-ind/full-ind/sections/} & The full-stage section files & Stage 3 sections; one \texttt{sec-NN-<name>.tex} per top-level section \\
\texttt{trial-ind/final-ind/sections/} & The final-stage section files & Stage 4 sections; the same one-to-one \texttt{sec-NN} mapping \\
\texttt{sec-06-cmc.tex} & Chemistry, Manufacturing and Control & File name matches the CMC section it contains \\
\texttt{sec-10-relevant-information.tex} & Relevant Information (this section) & File name matches the section it contains \\
\texttt{fig-09-five-grounding-verifications.md} & Source for Figure 21 & Figure file name matches the figure it renders \\
\bottomrule
\end{tabularx}
\figcaption{Table 11.3. Repository, directory, and file name-matching for the IND
(Verification 5). Every path named in the prose denotes the actual location on disk
in repository v4.3.0, making the package self-describing and navigable.}

The name-matching control has a second, quieter benefit for regulatory review:
completeness checking. Because the ReGARDD table of contents enumerates the
top-level sections, and because each top-level section corresponds to exactly one
\texttt{sec-NN-<name>.tex} file whose number and name match the contents entry, a
reviewer can verify that the package is complete by listing the
\texttt{sections/} directory and confirming that the numbered files run
consecutively with no gap. A missing section would appear as a missing file in the
sequence, and a misfiled section would appear as a name that does not match its
contents entry. The same logic applies to the figures: the 22 figure sources in
\texttt{trial-ind/mermaid/} are numbered consecutively, so a gap in the figure
numbering would be visible as a missing source file. This turns the naming
convention into a lightweight but real completeness check that does not depend on
reading the prose \cite{paisitedocpk, natlpaiplatf}.

\subsection{Limitations of the AI-Generation Method}

The grounding and verification controls above constrain the build but do not
eliminate its limitations, and those limitations are stated plainly here so that a
reviewer weighs the package accordingly. Table 11.4 enumerates the principal
limitations together with the mitigation applied to each. The most important point
is the one that no method-level control can change: the AI-generation method alters
none of the clinical realities of the trial. It does not change tumor biology, it
does not shorten the 28 day dose-limiting-toxicity observation window or the 30 and
90 day safety follow-up or the survival follow-up to 24 months, and it does not
compress the fixed 30 day FDA review clock or the IRB meeting calendar. What the
method compresses is the administrative and document-preparation bucket alone; the
clinical and operational bucket and the fixed regulatory-review bucket are
unchanged \cite{phase1guidance, pdac060s2030}. Every generated document, this IND
included, must still pass human medical, statistical, and regulatory review before
use, and the build is an independent research artifact, not an active IND or
investigational device exemption and not regulatory advice
\cite{cfr312paiuni, cfr050paiuni}.

Several narrower limitations follow from how the build operates. The method does not
rely on live internet search for a project of this size, because specifying single
external URLs is unreliable at scale; it depends instead on a curated repository of
appropriately sized inputs, which bounds the source material a single build can
absorb. Because the model does not always have access to a PDF compiler during
generation, a portion of senior-author formatting, including table column widths,
figure-caption spacing, and occasional page breaks, must still be applied by hand,
and the final stage performs that pass explicitly. Total input is kept under roughly
5 MB with individual files held under per-file token limits, which caps how much
source material a single build can ingest. High-dimensional or approximate inputs
can make outputs less deterministic, and oversized or overly complex inputs can
surface server errors that signal reduced quality; the build mitigates this by
keeping inputs curated and appropriately sized. These constraints are the same ones
characterized for the broader single-prompt method and are disclosed here for the
specific case of an IND \cite{ichpaistduni, chatgpt100kp, llmcomp16fda}.

A further limitation deserves explicit emphasis for a Phase 1 oncology submission.
The method's controls are controls on \textit{document construction}, not on
\textit{clinical judgment}. They ensure that the dose-escalation rule rendered in
the figure is the rule stated in the protocol, but they do not, and cannot, decide
whether 160 mg once daily is the correct starting dose, whether the 28 day DLT
window is the correct observation period, or whether the 3 N per-arm force cap is
the correct safety limit. Those decisions are made by the human investigators,
informed by the daraxonrasib clinical experience that supports the 300 mg once daily
RASolute 302 recommended Phase 2 dose, by the preclinical and computational evidence
base, and by the device engineering record, and they are subject to FDA and IRB
review. The AI-generation method is a faithful and auditable scribe of those
decisions; it is not a substitute for them. The boundary is the same staged-autonomy
boundary that governs the device: just as the on-premises language model in the
operating room is a second-opinion advisory oracle only, with full autonomy
prohibited under 21 CFR \S 312.21(e), the document-generation build is an advisory
authoring tool whose output is gated by human review \cite{congress9510,
clintrustpai}.

\needspace{12\baselineskip}
\noindent\begin{tabularx}{\textwidth}{L{4.8cm} Y}
\toprule
\textbf{Limitation} & \textbf{Mitigation applied in this IND} \\
\midrule
The method changes no clinical reality (tumor biology, DLT window, survival follow-up) & The 28 day DLT window, 30 and 90 day safety follow-up, and 24 month survival follow-up are preserved exactly as specified in the protocol; the build compresses only document preparation \\
The method cannot shorten fixed regulatory clocks & The 30 day FDA review and the IRB meeting calendar are stated as binding and unchanged; faster authoring only starts the clock sooner \\
The method controls construction, not clinical judgment & Dose levels, DLT window, and force caps are authored faithfully from the protocol but decided by human investigators under FDA and IRB review \\
Generated content can be fluent but wrong without grounding & The five name-matching and grounding verifications and the figure verify-twice checkpoint bind every figure and number to a named source file \\
No live internet search at scale; URLs unreliable in bulk & Inputs are a curated repository of appropriately sized files rather than ad hoc external URLs \\
No PDF compiler during generation; some formatting needs a manual pass & Senior-author formatting (column widths, caption spacing, page breaks) is applied explicitly in the final stage \\
Bounded input budget (under ~5 MB total; per-file token limits) & Source material is curated and apportioned so each build stays within limits without truncation \\
High-dimensional or oversized inputs reduce determinism and can surface server errors & Inputs are kept appropriately sized; complex content is split so the build remains deterministic and stable \\
The artifact is research, not an active regulatory submission & The package is labeled an independent research artifact, not an active IND or IDE and not regulatory advice; human review is required before any use \\
\bottomrule
\end{tabularx}
\vspace*{-0.3cm}
\figcaption{Table 11.4. Principal limitations of the AI-generation method and the
mitigation applied to each\\in this IND. The clinical and fixed-regulatory
limitations are intrinsic and cannot be removed; the\\method level limitations are
mitigated by grounding, verification, and a human review requirement.}

\subsubsection{A Worked Example of Grounding}

To make the grounding discipline concrete rather than abstract, consider a single
trial number followed from its source to its rendered occurrences. The recommended
Phase 2 dose of daraxonrasib, 300 mg once daily, is established in the daraxonrasib
clinical record (the RASolute 302 program) and is the value of the third and
highest dose level, DL3, in this trial's 3+3 escalation. From that one source the
number propagates, unchanged, to several places: it is DL3 in the dose-escalation
figure; it is the top of the dose-level ladder in the General Investigational Plan;
it is the dose at which a maximum tolerated dose determination would conclude if no
more than one of six subjects experienced a dose-limiting toxicity at that level;
and it is the dose carried into the precision discussion of the statistical plan,
where the maximum enrollment of 18 subjects (three levels of six) bounds the
confidence intervals. Table 11.5 traces this single value through its occurrences
and pairs it with the companion enrollment and precision figures, so a reviewer can
verify the entire chain from one source.

\needspace{10\baselineskip}
\noindent\begin{tabularx}{\textwidth}{L{4cm} L{3.4cm} Y}
\toprule
\textbf{Grounded value} & \textbf{Source of record} & \textbf{Where it recurs identically in the IND} \\
\midrule
300 mg once daily (DL3) & Daraxonrasib RASolute 302 RP2D & Top dose level in the dose-escalation figure, the General Investigational Plan ladder, and the MTD/RP2D narrative \\
28 day DLT window & Protocol v1.0.0 & Dose-escalation automaton, safety narrative, and the co-primary endpoint defining MTD/RP2D \\
18 treated subjects (3 levels $\times$ 6) & Protocol enrollment cap & Number-of-subjects figure, statistical precision table, and halt-rule denominators \\
Precision at 0 of 18 & Exact Clopper-Pearson & 0\% observed, 95\% CI 0\% to 18\%, stated identically in the statistical plan and any figure that cites it \\
Precision at 1 of 18 & Exact Clopper-Pearson & 5.6\% observed, 95\% CI 0.1\% to 27\%, stated identically wherever the single-event precision is discussed \\
\bottomrule
\end{tabularx}
\vspace*{-0.35cm}
\figcaption{Table 11.5. A worked grounding example. One source value, the 300 mg
once daily recommended Phase 2 dose, propagates unchanged through the
dose-escalation figure, the General Investigational Plan, and the statistical
precision discussion, alongside the 28 day DLT window, the 18 subject enrollment
cap, and the exact Clopper-Pearson precision figures, so a reviewer can verify the
whole chain from a single source.}

The worth of the example is that it shows grounding operating at the level of
review, not just of authoring. A reviewer who confirms that DL3 is 300 mg once
daily in the dose-escalation figure, and that the source of record for that value
is the daraxonrasib RASolute 302 recommended Phase 2 dose, has simultaneously
confirmed every other place the value appears, because the grounding discipline
guarantees they are the same value drawn from the same source. The same holds for
the 28 day DLT window, the 18 subject cap, and the exact Clopper-Pearson precision
figures of 0 percent (95 percent confidence interval 0 percent to 18 percent) at 0
of 18 and 5.6 percent (95 percent confidence interval 0.1 percent to 27 percent) at
1 of 18. The grounding does not make these clinical or statistical choices correct;
it makes them consistent and auditable, which is exactly the property a reviewer
needs to evaluate them efficiently \cite{oncotrialllm, pdac060s2030}.

\subsection{Compilation Statement}

All LaTeX sources for this IND compile in Overleaf with pdfLaTeX. The final-stage
document is built from \texttt{main.tex}, the shared style file
\texttt{indstyle.sty}, the bibliography \texttt{references.bib}, and the
\texttt{sections/} directory of section files, with every figure rendered as vector
TikZ inside the \texttt{mermaidfig} environment and every table set full width with
the ragged-right column types defined in the style file. There are no raster images,
no \texttt{\textbackslash includegraphics} calls, and no external image
dependencies, so the package compiles from its own sources without additional
assets. The bibliography is processed with the \texttt{ieeetr} style, and the
codified references throughout the package use the section symbol, for example 21
CFR \S 312.23, exactly as rendered by the style file. A reviewer can therefore
reproduce the compiled PDF directly from the repository at version v4.3.0 by
uploading the stage directory to Overleaf and compiling with pdfLaTeX, which closes
the loop from grounded source to verifiable, reproducible document. The
reproducibility is exact in the sense that the same sources produce the same
document: the figures are deterministic TikZ with fixed coordinates, the tables are
fixed-width tabularx with explicitly stated column widths, and the style file pins
the grayscale ramp, the page geometry, and the typographic penalties, so a
third-party compile yields a PDF that matches the submitted one figure for figure
and table for table \cite{phase1guidance, daraxwhipple, onpremwhippl}.
